\chapter{Verificación y Depuración del Sistema}
\label{ch:resultados}

\section{Objetivos de las Pruebas}
\label{sec:objetivos_pruebas}

El objetivo general del proceso de verificación fue confirmar que la
arquitectura HR-RT-Reporter, en su estado final, cumple de punta a punta
con la propuesta metodológica de este trabajo: capturar eventos de
ejecución del núcleo Ibex con efecto sonda nulo, fecharlos con una base de
tiempo de hardware, y ponerlos a disposición del sistema operativo Linux
de forma confiable, sin pérdida de eventos por metaestabilidad entre
dominios de reloj. De este objetivo general se desprenden los siguientes
objetivos específicos, cada uno abordado en una sección independiente de
este capítulo:

\begin{itemize}
    \item Verificar que el canal de comunicación entre el Procesador de
    Sistema (PS) y la Lógica Programable (PL) -- disparo de arranque por
    GPIO, acceso a la BRAM, exclusión mutua -- funciona como fue diseñado
    (Sección~\ref{sec:pruebas_ps_pl}).
    \item Verificar que el núcleo Ibex ejecuta el firmware correctamente
    desde la memoria de instrucciones local, sin quedar en estados
    indeterminados ni ejecutar instrucciones incorrectas tras un salto
    (Sección~\ref{sec:pruebas_ibex}).
    \item Verificar que el módulo Reporter captura los eventos de
    escritura del Ibex en las direcciones correctas
    (Sección~\ref{sec:pruebas_eventos}).
    \item Verificar que la base de tiempo de hardware (\textit{timestamp})
    es monótona, independiente y correctamente accesible
    (Sección~\ref{sec:verificacion_timestamps}).
    \item Verificar que el cruce entre los dos dominios de reloj
    asíncronos del sistema (Ibex a $50\text{ MHz}$, Reporter a
    $100\text{ MHz}$) no introduce pérdida de eventos
    (Sección~\ref{sec:verificacion_cdc}).
    \item Verificar que el sistema operativo Linux puede leer de forma
    confiable, vía AXI-Lite, los resultados capturados por el hardware
    (Sección~\ref{sec:pruebas_linux}).
\end{itemize}

Un objetivo transversal a todas las pruebas anteriores fue, siempre que
resultara posible, validarlas primero mediante simulación funcional
(\texttt{xsim}) antes de comprometer un ciclo completo de síntesis,
implementación y generación de \textit{bitstream} -- del orden de $15$ a
$20$ minutos por corrida -- para probarlas sobre la placa física.

%--------------------------------------------------------------------------------
%5.2) METODOLOGIA DE VERIFICACION
%--------------------------------------------------------------------------------

\section{Metodología de Verificación}
\label{sec:metodologia_debug}

\subsection{Preparación del Entorno de Simulación}
\label{ssec:entorno_xsim}

Antes de poder ejecutar cualquier simulación de línea de comandos con las
herramientas de Vivado (\texttt{xvlog}, \texttt{xelab}, \texttt{xsim}), es
necesario preparar el entorno de la siguiente manera:

\begin{enumerate}
    \item \textbf{Habilitar las herramientas en el \texttt{PATH}:} abrir
    una consola (símbolo del sistema o PowerShell) y ejecutar el script de
    configuración de la instalación de Vivado:
    
\begin{lstlisting}[style=powershellstyle]
call E:\vv\Vivado\2022.2\settings64.bat
\end{lstlisting}

    Este script agrega al \texttt{PATH} de la sesión actual el directorio
    \texttt{bin} de Vivado (donde residen \texttt{xvlog.exe},
    \texttt{xelab.exe} y \texttt{xsim.exe}) y define las variables de
    entorno necesarias (\texttt{XILINX\_VIVADO}, licencias, etc.). Es el
    mismo entorno que usa Vivado internamente para su función
    \textit{"Run Simulation"} desde la interfaz gráfica, por lo que los
    resultados obtenidos por línea de comandos son equivalentes a los que
    se obtendrían desde la GUI.
    \item \textbf{Ubicarse en el directorio del proyecto:} los comandos de
    simulación se ejecutan desde el directorio raíz del proyecto
    (\texttt{C:\textbackslash pro\textbackslash reporter5}), de forma que
    las rutas relativas a los archivos fuente del \textit{fileset} de
    simulación (\texttt{reporter5.srcs/sim\_1/new/}) sean consistentes.
    \item \textbf{Verificar que el entorno quedó configurado:} ejecutar
    \texttt{xvlog -version} y confirmar que reporta la versión de Vivado
    instalada (\texttt{2022.2} en este proyecto). Si el comando no se
    reconoce, el paso 1 no se ejecutó correctamente en esa misma consola.
\end{enumerate}

%--------------------------------------------------------------------------------
%5.2) AGREGAR EL TESTBENCH AL PROYECTO Y CONFIGURARLO COMO TOP
%--------------------------------------------------------------------------------

\subsection{Incorporación del \textit{Testbench} y Configuración como Módulo Superior}
\label{ssec:tb_como_top}

Con el entorno preparado, el siguiente paso es incorporar el archivo de
banco de pruebas (\textit{testbench}) al conjunto de fuentes de
simulación del proyecto de Vivado (\texttt{reporter5.xpr}):

\begin{enumerate}
    \item \textbf{Agregar el archivo al \textit{fileset} de simulación:}
    en la ventana \textit{Sources} de Vivado, clic derecho sobre
    \textit{Simulation Sources (sim\_1)} $\rightarrow$ \textit{Add
    Sources...} $\rightarrow$ \textit{Add or create simulation sources},
    incorporando el archivo \texttt{tb\_ibex\_bram\_subsystem.sv}
    (ubicado en \texttt{reporter5.srcs/sim\_1/new/}). El mismo resultado
    se obtiene por línea de comandos Tcl con:
    
\begin{lstlisting}[style=bashstyle]
add_files -fileset sim_1 reporter5.srcs/sim_1/new/tb_ibex_bram_subsystem.sv
\end{lstlisting}

    \item \textbf{Configurar el \textit{testbench} como módulo superior
    (\textit{Set as Top}):} un \textit{fileset} de simulación puede
    contener más de un banco de pruebas (en este proyecto conviven, por
    ejemplo, distintas variantes de instrumentación del mismo archivo a lo
    largo de las sesiones de depuración), por lo que es necesario indicarle
    a Vivado cuál debe instanciarse como raíz de la jerarquía de
    simulación. En la GUI: clic derecho sobre el módulo
    \texttt{tb\_ibex\_bram\_subsystem} en el árbol de fuentes
    $\rightarrow$ \textit{Set as Top}. Por línea de comandos, la propiedad
    equivalente es:

\begin{lstlisting}[style=bashstyle]
set_property top tb_ibex_bram_subsystem [get_filesets sim_1]
set_property top_lib xil_defaultlib [get_filesets sim_1]
\end{lstlisting}

    \item \textbf{Dependencias resueltas automáticamente:} a partir de este
    punto, Vivado (o \texttt{xvlog}/\texttt{xelab} invocados manualmente,
    ver Sección~\ref{ssec:ejecucion_xsim}) resuelve el orden de análisis de
    todos los archivos fuente de los que depende el \textit{testbench}: los
    paquetes y módulos del núcleo Ibex (importados vía FuseSoC, ver
    Sección~\ref{ssec:fuse}), el subsistema propio
    \texttt{ibex\_bram\_subsystem.sv} y, en la variante extendida del
    \textit{testbench} usada para validar el mecanismo de cruce de
    dominios de reloj, el IP completo del Reporter
    (\texttt{axi\_reporter\_v1\_0}, con sus tres archivos fuente:
    \texttt{reporter\_core.sv}, \texttt{axi\_reporter\_v1\_0\_S00\_AXI.v}
    y \texttt{axi\_reporter\_v1\_0.v}).
\end{enumerate}
El código fuente completo de esta variante extendida del \textit{testbench} -- incluyendo la instrumentación de traza, el \textit{driver} AXI-Lite y las verificaciones automáticas contra los
    valores esperados -- se incluye íntegro en el Apéndice~\ref{ap:script-tb}.
%--------------------------------------------------------------------------------
%5.3) EJECUCION DE LA SIMULACION (xsim)
%--------------------------------------------------------------------------------

\subsection{Ejecución de la Simulación}
\label{ssec:ejecucion_xsim}

Con el \textit{testbench} configurado como módulo superior, existen dos
vías equivalentes para ejecutarlo, ambas invocando internamente las mismas
tres herramientas:

\begin{itemize}
    \item \textbf{Desde la interfaz gráfica de Vivado:} pulsando
    \textit{Run Simulation} $\rightarrow$ \textit{Run Behavioral
    Simulation}, Vivado compila, elabora y lanza automáticamente el visor
    de formas de onda (\textit{waveform viewer}), útil para una inspección
    visual rápida.
    \item \textbf{Desde la línea de comandos:} la vía utilizada
    preferentemente durante el ciclo intensivo de depuración de este
    trabajo, ya que permite iterar con mayor velocidad (sin la sobrecarga
    de reabrir la interfaz gráfica en cada corrida) y automatizar la
    búsqueda de patrones de error en los registros de salida
    (\textit{logs}) mediante herramientas de línea de comandos. El flujo
    consta de tres invocaciones sucesivas:
    \begin{enumerate}
        \item \texttt{xvlog} compila cada archivo fuente en el orden de
        dependencias correcto, hacia una biblioteca de simulación
        (\texttt{xil\_defaultlib} por defecto).
        \item \texttt{xelab} elabora la jerarquía completa del diseño a
        partir del módulo superior configurado, y genera un
        \textit{snapshot} ejecutable de la simulación, por ejemplo:
\begin{verbatim}
xelab --incr --debug typical --relax --mt 2 -L xil_defaultlib
  -L blk_mem_gen_v8_4_5 -L uvm -L xilinx_vip -L unisims_ver
  -L unimacro_ver -L secureip -L xpm --snapshot tb_reporter_cdc_behav
  xil_defaultlib.tb_ibex_bram_subsystem xil_defaultlib.glbl -log xelab.log
\end{verbatim}
        (se incluye siempre \texttt{xil\_defaultlib.glbl}, el módulo
        global de señales de Xilinx requerido por las primitivas
        \textit{unisim} como \texttt{BUFGCE} y por \texttt{blk\_mem\_gen}).
        \item \texttt{xsim} ejecuta ese \textit{snapshot} de principio a
        fin y vuelca la salida de las sentencias \texttt{\$display} del
        \textit{testbench} a un archivo de registro para su análisis
        posterior:
\begin{verbatim}
xsim tb_reporter_cdc_behav -log xsim_run.log -runall
\end{verbatim}
    \end{enumerate}
\end{itemize}

%--------------------------------------------------------------------------------
%5.4) METODOLOGIA DE DIAGNOSTICO: COMO SE LLEGO A DESCIFRAR CADA ANOMALIA
%--------------------------------------------------------------------------------

\subsection{Metodología de Diagnóstico: Cómo se Identificó Cada Anomalía}
\label{ssec:metodologia_diagnostico}

Una vez operativo el flujo de simulación, la identificación de cada
defecto siguió un proceso sistemático de refinamiento progresivo, en el
que cada corrida exitosa habilitaba la detección del problema siguiente,
inicialmente enmascarado por el anterior:

\begin{enumerate}
    \item \textbf{Primer filtro -- errores de compilación/elaboración:} la
    primera corrida de \texttt{xelab} fallaba con un error de compilación
    del código C++ generado internamente por la herramienta a partir de
    una exportación \texttt{DPI-C} dentro de un bloque \texttt{generate}
    con nombre del RTL de Ibex. Esto impedía incluso obtener un
    \textit{snapshot} ejecutable, y se resolvió aislando esa exportación
    bajo una directiva de compilación condicional (\texttt{`ifdef
    VERILATOR}) que nunca está activa en el flujo de Vivado.

    \item \textbf{Segundo filtro -- aserciones internas del núcleo:} con la
    simulación ya elaborando correctamente, el registro de salida mostraba
    decenas de fallos de aserción por segundo de simulación
    (\texttt{ASSERT FAILED}) desde el segundo ciclo en adelante, causados
    por puertos de \texttt{ibex\_top} sin conectar que permanecían en
    estado indeterminado (\texttt{'X'}). Se identificaron mediante una
    búsqueda de la cadena \texttt{"ASSERT FAILED"} en el registro,
    cruzada contra la lista de puertos de entrada del núcleo sin valor por
    defecto.

    \item \textbf{Tercer filtro -- instrumentación de traza ciclo a
    ciclo:} superadas las dos anomalías anteriores, la simulación se
    comportaba correctamente durante el arranque, pero el núcleo caía en
    un bucle de excepción permanente inmediatamente después del primer
    salto tomado del programa. Para diagnosticar la causa raíz, se
    instrumentó el \textit{testbench} con sentencias \texttt{\$display}
    disparadas en cada flanco de reloj, volcando el estado de las señales
    de \textit{fetch} de instrucciones y de acceso a datos (dirección,
    dato, señal de válido) durante los primeros $300$ ciclos posteriores a
    la liberación del reset. Esta traza se contrastó manualmente, ciclo
    por ciclo, contra el contenido real del archivo de inicialización
    (\texttt{.mif}/\texttt{.coe}) de la memoria de instrucciones --
    decodificado a mano, con el desplazamiento de dirección correcto --
    lo que permitió observar que el dato capturado inmediatamente después
    de cada salto correspondía a la dirección \textit{anterior}, todavía
    en tránsito dentro de la memoria, revelando un ciclo de latencia
    faltante en el modelo de la señal de válido de lectura frente a la
    latencia real de dos ciclos de la primitiva de memoria instanciada
    (ver Sección~\ref{ssec:bram_dual_port}).

    \item \textbf{Cuarto filtro -- verificación positiva de extremo a
    extremo:} agotados los defectos detectables por fallos evidentes en el
    registro de simulación, el \textit{testbench} se extendió para
    instanciar además el periférico Reporter completo, con su propio
    generador de reloj independiente de $100\text{ MHz}$ (asíncrono
    respecto del reloj de $50\text{ MHz}$ del Ibex), junto con un
    \textit{driver} AXI4-Lite de lectura escrito a mano. Esto permitió,
    por primera vez, contrastar explícitamente los valores leídos por
    simulación contra los valores esperados de forma analítica (por
    ejemplo, el resultado final de la secuencia pseudoaleatoria generada
    por el firmware), confirmando -- o, en un caso, refutando -- que cada
    corrección aplicada al mecanismo de captura y sincronización se
    comportaba como se esperaba, antes de invertir un ciclo completo de
    síntesis e implementación para probarla en la placa física.
\end{enumerate}

Este proceso de refinamiento progresivo -- en el que cada capa de defectos
ocultaba a la siguiente -- ilustra por qué la verificación por simulación
resultó indispensable en este trabajo: varios de estos problemas, en
particular el de latencia de memoria del tercer filtro, habrían sido
extremadamente difíciles de diagnosticar directamente sobre hardware
físico, donde no existe visibilidad ciclo a ciclo de las señales internas
sin instrumentación adicional (por ejemplo, un analizador lógico integrado
o \textit{ILA}).

%--------------------------------------------------------------------------------
%5.3) PRUEBAS DE COMUNICACION PS-PL
%--------------------------------------------------------------------------------

\section{Pruebas de Comunicación PS--PL}
\label{sec:pruebas_ps_pl}

Antes de ejercitar el núcleo Ibex o el módulo Reporter en sí, se verificó
que los tres canales de comunicación entre el Procesador de Sistema y la
Lógica Programable funcionan correctamente:

\begin{itemize}
    \item \textbf{Disparo de arranque por GPIO:} el script
    \texttt{recolector.py} (Sección~\ref{sec:pruebas_linux}) controla la
    línea de reinicio del Ibex a través de \texttt{axi\_gpio\_0}
    (\texttt{0x41200000}), escribiendo primero un $0$ para mantener el
    núcleo en reposo y luego un $1$ para liberarlo. Se verificó, tanto en
    simulación (manteniendo \texttt{rst\_ni} en bajo) como en hardware
    real, que sin esta liberación explícita del reset el núcleo Ibex
    permanece detenido y ningún evento llega a capturarse.
    \item \textbf{Acceso a la memoria BRAM vía
    \texttt{axi\_bram\_ctrl\_0}:} se verificó que, con el Ibex retenido en
    reset, el ARM puede inicializar y leer la memoria local a través de la
    ventana de $8\text{ KB}$ expuesta en \texttt{0x40000000} (ver
    Sección~\ref{sssec:memory_map} para la distinción entre esta ventana y
    la profundidad real de $16\text{ KiB}$ de la BRAM).
    \item \textbf{Exclusión mutua:} se verificó que, una vez liberado el
    reset, el multiplexor de exclusión mutua
    (Sección~\ref{ssec:exclusion_mutua}) revoca los privilegios de
    escritura del ARM sobre la BRAM, protegiendo la integridad del
    firmware en ejecución frente a accesos concurrentes desde Linux.
\end{itemize}

%--------------------------------------------------------------------------------
%5.4) PRUEBAS DEL NUCLEO IBEX
%--------------------------------------------------------------------------------

\section{Pruebas del Núcleo Ibex}
\label{sec:pruebas_ibex}

Las primeras corridas de simulación se concentraron exclusivamente en el
subsistema Ibex+BRAM (sin el periférico Reporter), con el objetivo de
confirmar que el núcleo ejecuta el firmware correctamente antes de
involucrar al resto de la arquitectura. Este proceso -- descripto en
detalle como el primer, segundo y tercer filtro de la
Sección~\ref{ssec:metodologia_diagnostico} -- permitió identificar y
corregir tres defectos: un fallo de compilación del código DPI generado
por la herramienta a partir del RTL de Ibex, decenas de fallos de
aserción por puertos del núcleo sin conectar, y un desajuste de un ciclo
en el modelo de latencia de lectura de la BRAM que provocaba una
instrucción ilegal inmediatamente después de cualquier salto tomado. Una
vez corregidos los tres, la simulación mostró al núcleo ejecutando el
firmware completo sin fallos de aserción y sin instrucciones ilegales
durante los $20\,\mu s$ simulados, incluyendo el comportamiento esperado
de \textit{fetch} especulativo alrededor del bucle infinito final del
programa.

%--------------------------------------------------------------------------------
%5.5) PRUEBAS DE GENERACION DE EVENTOS
%--------------------------------------------------------------------------------

\section{Pruebas de Generación de Eventos}
\label{sec:pruebas_eventos}

Con el núcleo Ibex verificado, el siguiente objetivo fue confirmar que el
módulo Reporter captura correctamente los eventos que el firmware genera
mediante escrituras a las direcciones reservadas (\texttt{0x80000008},
\texttt{0x8000000C} y \texttt{0x80000010}, Sección~\ref{sec:axi_reporter}).
Esta verificación expuso dos defectos relevantes: primero, que la lógica
de captura comparaba contra una única dirección (\texttt{0x80000000}) que
el firmware nunca utiliza -- por lo que el contador de diagnóstico daba
siempre cero en hardware, de forma perfectamente reproducible --,
corregido comparando contra las direcciones reales; y segundo, que el mapa
de registros de lectura por AXI-Lite no alcanzaba a exponer todos los
eventos capturados (limitado originalmente por un ancho de dirección
insuficiente), ampliado a los $8$ registros de la
Tabla~\ref{tab:mapa_registros_reporter}. Tras ambas correcciones, la
simulación confirmó que cada escritura del Ibex a una de las tres
direcciones reservadas se refleja correctamente en el registro
\textit{sticky} correspondiente (\texttt{EVT\_PC}, \texttt{EVT\_I} o
\texttt{EVT\_D}).

%--------------------------------------------------------------------------------
%5.6) VERIFICACION DE LA ADQUISICION DE TIMESTAMPS
%--------------------------------------------------------------------------------

\section{Verificación de la Adquisición de \textit{Timestamps}}
\label{sec:verificacion_timestamps}

El contador libre de $64$ bits que provee la base de tiempo de hardware
(\texttt{TIMESTAMP\_LO}/\texttt{TIMESTAMP\_HI}, Tabla~
\ref{tab:mapa_registros_reporter}) se verificó de forma independiente del
resto del mecanismo de captura de eventos: al operar en el dominio de
$100\text{ MHz}$ propio del Reporter, se confirmó en simulación que se
incrementa de forma monótona y continua desde el instante en que se libera
el reset general del sistema, sin depender de que el núcleo Ibex esté o no
ejecutando firmware. Esta independencia es intencional: permite que Linux
feche sus propias lecturas periódicas del Reporter con una base de tiempo
común, sin necesidad de que el Ibex participe en absoluto de la generación
de la marca temporal.

%--------------------------------------------------------------------------------
%5.7) VERIFICACION DE LA SINCRONIZACION ENTRE DOMINIOS DE RELOJ
%--------------------------------------------------------------------------------

\section{Verificación de la Sincronización entre Dominios de Reloj}
\label{sec:verificacion_cdc}

Esta es la prueba más crítica de todo el proceso de verificación, dado que
el problema de cruce de dominios de reloj (\textit{Clock Domain Crossing},
CDC) fue, en una primera instancia, la causa raíz del síntoma reportado
originalmente: un contador de diagnóstico que permanecía en cero de forma
reproducible. Para poder reproducir y verificar este mecanismo en
simulación -- y no solo por inspección de código -- el \textit{testbench}
se extendió para instanciar el periférico Reporter completo con su propio
generador de reloj de $100\text{ MHz}$, independiente y asíncrono respecto
del reloj de $50\text{ MHz}$ del Ibex (Sección~\ref{ssec:tb_como_top}),
junto con un \textit{driver} AXI-Lite de lectura que permite contrastar
explícitamente los contadores de diagnóstico \texttt{DBG\_RAW\_COUNT}
(flancos de escritura sin sincronizar, expuestos deliberadamente sin
protección para fines de comparación) y \texttt{DBG\_SYNC\_COUNT} (flancos
ya resueltos por la cadena sincronizadora de tres \textit{flip-flops}). En
operación correcta, ambos contadores deben coincidir exactamente, ya que
cualquier divergencia entre ellos indicaría pérdida de eventos por
metaestabilidad. La Tabla~\ref{tab:resultados_sim}
(Sección~\ref{ssec:resultados_simulacion}) muestra que, en todas las
corridas realizadas -- con firmwares de $2$ y de $3000$ escrituras --,
\texttt{DBG\_RAW\_COUNT} y \texttt{DBG\_SYNC\_COUNT} coincidieron
exactamente, confirmando que el sincronizador no pierde eventos bajo las
condiciones de carga ensayadas.

%--------------------------------------------------------------------------------
%5.8) PRUEBAS DE LECTURA DESDE LINUX
%--------------------------------------------------------------------------------

\section{Pruebas de Lectura desde Linux}
\label{sec:pruebas_linux}

La última etapa del proceso de verificación consistió en confirmar que el
sistema operativo Linux, ejecutándose en el Procesador de Sistema, puede
leer de forma confiable los resultados capturados por el hardware. Esta
prueba se realizó en dos niveles: primero, replicando en el
\textit{testbench} de simulación (Sección~\ref{ssec:tb_como_top}) la misma
secuencia de lecturas AXI-Lite que realiza el script
\texttt{recolector.py} sobre el hardware real (disparo del reset por GPIO,
seguido de una secuencia de lecturas sobre los $8$ registros del
Reporter); y segundo, ejecutando el propio \texttt{recolector.py} sobre la
placa Zybo reprogramada, guardando el resultado en
\texttt{trazabilidad\_ibex.txt} para su comparación posterior
(Sección~\ref{ssec:resultados_hardware}). Esta prueba también dejó una
lección metodológica relevante: durante el desarrollo, el mapa de
registros del Reporter creció incorporando el registro \texttt{EVT\_I} en
un \textit{offset} previamente sin uso, lo que obligó a actualizar en
consecuencia los \textit{offsets} leídos por \texttt{recolector.py} -- una
corrida con el script desactualizado produce columnas con la etiqueta
equivocada (aunque los datos subyacentes sigan siendo correctos),
ilustrando por qué conviene verificar cada \textit{offset} contra el RTL
vigente antes de interpretar los resultados de una corrida.

%--------------------------------------------------------------------------------
%5.9) PROBLEMAS DETECTADOS Y CORRECCIONES IMPLEMENTADAS
%--------------------------------------------------------------------------------

\section{Problemas Detectados y Correcciones Implementadas}
\label{sec:problemas_correcciones}

A lo largo de las pruebas descriptas en las secciones anteriores se
identificaron y corrigieron diez defectos independientes, resumidos en la
Tabla~\ref{tab:resumen_bugs}.

\begin{table}[htbp]
\centering
\small
\caption{Resumen de los defectos detectados y corregidos durante el proceso de verificación.}
\label{tab:resumen_bugs}
\begin{tabular}{p{0.5cm}p{3.6cm}p{4.3cm}p{4.3cm}}
\hline
\textbf{\#} & \textbf{Síntoma} & \textbf{Causa raíz} & \textbf{Corrección} \\ \hline
1 & Fallo de compilación en \texttt{xelab} & Exportación \texttt{DPI-C} dentro de un bloque \texttt{generate} con nombre & Exportación aislada bajo \texttt{`ifdef VERILATOR} \\
2 & Decenas de \texttt{ASSERT FAILED} por segundo de simulación & Puertos de \texttt{ibex\_top} sin conectar, en estado \texttt{'X'} & Conexión de esos puertos a constantes seguras \\
3 & Instrucción ilegal inmediatamente después del primer salto tomado & Latencia de la BRAM modelada en 1 ciclo (real: 2 ciclos) & Cadena de 2 \textit{flip-flops} para la señal de válido de lectura \\
4 & Contador de diagnóstico siempre en 0 en hardware & Captura de eventos comparaba contra una única dirección que el firmware nunca usa & Comparación contra las direcciones reales de escritura \\
5 & El fix del \#4 no llegaba al \textit{bitstream} regenerado & \texttt{.xci} del IP apuntando a un directorio de generación residual de una sesión de \textit{IP Packager} & Replicación manual del fix en ambas copias (pendiente resolver de fondo) \\
6--7 & \textit{Offsets} de diagnóstico aliaseados; sin sincronizador de CDC & Mapa de registros incompleto; muestreo combinacional entre dominios de reloj & Mapa de 8 registros; registros \textit{sticky} + sincronizador de 3 \textit{flip-flops} \\
8 & Los fixes de un IP no se reflejaban tras \texttt{reset\_run} & Las corridas de síntesis \textit{Out-of-Context} de cada IP no se reseteaban junto con la del proyecto top & \texttt{reset\_run} explícito de todas las corridas, incluidas las OOC \\
9 & \textit{Driver} AXI-Lite del \textit{testbench} colgado & Carrera de muestreo en el mismo flanco de \texttt{posedge} & Muestreo en \texttt{negedge}, \texttt{ARVALID} sostenido hasta \texttt{RVALID} \\
10 & \texttt{EVT\_I} inexistente pese a ser asumido por firmware y script & Registro nunca implementado en el RTL & Registro \texttt{evt\_i\_o} agregado en el offset \texttt{0x14} \\ \hline
\end{tabular}
\end{table}


%--------------------------------------------------------------------------------
%5.10) RESULTADOS EXPERIMENTALES
%--------------------------------------------------------------------------------

\section{Resultados Experimentales}
\label{sec:resultados}

\subsection{Resultados de Simulación de Extremo a Extremo}
\label{ssec:resultados_simulacion}

La Tabla~\ref{tab:resultados_sim} resume los valores obtenidos en las
corridas de simulación con el \textit{testbench} extendido (\textit{driver}
AXI-Lite leyendo el mapa de registros completo del Reporter,
Tabla~\ref{tab:mapa_registros_reporter}), para dos firmwares distintos, en
contraste con los valores esperados calculados analíticamente:

\begin{table}[htbp]
\centering
\caption{Resultados de la simulación de extremo a extremo (\textit{snapshot} \texttt{tb\_reporter\_cdc\_behav}), valores esperados vs. obtenidos.}
\label{tab:resultados_sim}
\begin{tabular}{lll}
\hline
\textbf{Firmware} & \textbf{Registro} & \textbf{Valor esperado / obtenido} \\ \hline
\texttt{sanity\_check.c} & \texttt{DBG\_RAW\_COUNT} / \texttt{DBG\_SYNC\_COUNT} & $2$ / $2$ \\
\texttt{sanity\_check.c} & \texttt{EVT\_PC} & \texttt{0xCAFEF00D} \\
\texttt{sanity\_check.c} & \texttt{EVT\_D} & \texttt{0xDEADBEEF} \\
\texttt{reporter.c} ($N=1000$) & \texttt{DBG\_RAW\_COUNT} / \texttt{DBG\_SYNC\_COUNT} & $3000$ / $3000$ \\
\texttt{reporter.c} ($N=1000$) & \texttt{EVT\_PC} & \texttt{0x0000010C} \\
\texttt{reporter.c} ($N=1000$) & \texttt{EVT\_I} & $999$ \\
\texttt{reporter.c} ($N=1000$) & \texttt{EVT\_D} & \texttt{0x14E875E1} \\ \hline
\end{tabular}
\end{table}

En ambos casos, la totalidad de los valores leídos por el \textit{driver}
AXI-Lite coincidió exactamente con el valor esperado, validando en
simulación el mecanismo completo de captura y sincronización entre
dominios de reloj antes de proceder a su verificación en hardware real.

%--------------------------------------------------------------------------------
%(sigue dentro de 5.10, como subsection) RESULTADOS EN HARDWARE REAL
%--------------------------------------------------------------------------------

\subsection{Resultados en Hardware Real}
\label{ssec:resultados_hardware}

Con el bitstream generado a partir del diseño verificado en simulación, se
reprogramó la placa Zybo y se ejecutó \texttt{recolector.py} sobre el
sistema Linux embebido, con el firmware \texttt{reporter.c} ($N=1000$)
precargado en la BRAM de instrucciones. De las muestras válidas
registradas en \texttt{trazabilidad\_ibex.txt}, se obtuvo:
\begin{itemize}
    \item \texttt{EVT\_PC} constante en \texttt{0x0000010C} en todas las
    muestras -- coincide exacto con el valor predicho por simulación.
    \item Valor final de \texttt{EVT\_D}: $350{,}778{,}849$ decimal, es
    decir \texttt{0x14E875E1} en hexadecimal -- coincide exacto con el
    valor predicho por simulación.
    \item Valor final de \texttt{EVT\_I}: $999$ -- coincide exacto con el
    valor predicho por simulación.
\end{itemize}

Los tres valores, observables de forma completamente independiente entre
sí, coinciden de manera exacta entre lo predicho analíticamente, lo
verificado en simulación y lo medido en la placa física real, validando de
punta a punta la cadena completa: núcleo Ibex ($50\text{ MHz}$) $\rightarrow$
memoria de instrucciones $\rightarrow$ captura y sincronización de eventos
en \texttt{reporter\_core.sv} ($\text{CDC } 50 \rightarrow 100\text{ MHz}$)
$\rightarrow$ banco de registros AXI-Lite $\rightarrow$ recolección desde
Linux ($\texttt{recolector.py}$).

\textbf{Nota metodológica:} durante el desarrollo, el mapa de registros del
Reporter creció incorporando el registro \texttt{EVT\_I} (offset
\texttt{0x14}, previamente sin uso), lo que obligó a actualizar en
consecuencia los \textit{offsets} leídos por \texttt{recolector.py}. Este
episodio ilustra por qué, en un diseño que evoluciona, verificar cada
\textit{offset} contra el RTL vigente -- y no solo contra la documentación
o versiones anteriores del script de recolección -- resulta indispensable
antes de interpretar los resultados de una corrida.

%--------------------------------------------------------------------------------
%5.11) ANALISIS DE RESULTADOS
%--------------------------------------------------------------------------------

\section{Análisis de Resultados}
\label{sec:analisis_resultados}

La coincidencia exacta, en las tres instancias observadas de forma
independiente (\texttt{EVT\_PC}, \texttt{EVT\_D} y \texttt{EVT\_I}), entre
los valores predichos analíticamente, los obtenidos en simulación y los
medidos en la placa física real, constituye evidencia sólida de que la
arquitectura funciona de punta a punta tal como fue diseñada, y de que el
proceso de verificación por simulación descripto en la
Sección~\ref{sec:metodologia_debug} es representativo del comportamiento
real del hardware sintetizado. Asimismo, la coincidencia entre
\texttt{DBG\_RAW\_COUNT} y \texttt{DBG\_SYNC\_COUNT} en todas las corridas
realizadas (Sección~\ref{sec:verificacion_cdc}) indica que el
sincronizador de cruce de dominios de reloj no pierde eventos bajo las
condiciones de carga ensayadas hasta el momento (hasta $3000$ escrituras
consecutivas, con el firmware \texttt{reporter.c}).

No obstante, corresponde señalar explícitamente los límites de lo
verificado hasta el momento:

\begin{itemize}
    \item El módulo Reporter retiene únicamente el \textbf{último} evento
    capturado en cada registro \textit{sticky}, y no mantiene un historial
    ni una cola (\textit{FIFO}) de eventos previos. En consecuencia,
    cuando el firmware ejecuta muchas más escrituras de las que el script
    de recolección alcanza a muestrear en el mismo intervalo de tiempo, la
    mayoría de las lecturas observan el mismo valor final, no una traza
    continua de la ejecución -- una limitación de la arquitectura, no un
    defecto de la implementación.
    \item El camino de \textbf{lectura} del Reporter desde el lado del
    Ibex (por ejemplo, una eventual instrucción \texttt{lw} que consultara
    \texttt{current\_timestamp} desde el firmware) no cuenta con el mismo
    tratamiento de sincronización que el camino de escritura, y no fue
    ejercitado por ningún firmware utilizado hasta el momento (ni
    \texttt{sanity\_check.c} ni \texttt{reporter.c} leen del Reporter). Si
    en el futuro se incorpora firmware que realice esta operación, este
    camino requeriría el mismo análisis de CDC aplicado al camino de
    escritura antes de considerarse verificado.
    \item Las pruebas de sincronización entre dominios de reloj se
    realizaron bajo la carga de escritura generada por los dos firmwares
    disponibles (hasta $3000$ escrituras consecutivas a $50\text{ MHz}$);
    no se ensayaron patrones de escritura más agresivos (por ejemplo, una
    escritura por ciclo de forma sostenida) que pudieran ejercitar de
    manera más exigente al sincronizador.
\end{itemize}

En conjunto, los resultados obtenidos permiten concluir que la hipótesis
central de este trabajo -- la viabilidad de un mecanismo de trazabilidad
de hardware con efecto sonda nulo, robusto frente al cruce entre dominios
de reloj asíncronos -- queda validada para el alcance efectivamente
ensayado, dejando explícitamente identificados los aspectos que excedieron
ese alcance como trabajo futuro.